\chapter{Evaluierungskonzept und Ergebnisse}


\section{Evaluierung von Retrieval- und Ranking-Systemen}

\subsection{Recall@K}

Recall gehört zu den grundlegenden Metriken des Information Retrieval \cite{manning2008ir}.

\begin{align}
  Recall@K = \frac{ \# \text{ relevante Treffer in Top-K} } { \# \text{ relevante Treffer} }
\end{align}

Im vorliegenden System werden insbesondere

\begin{itemize}
  \item target\_recall@1
  \item target\_recall@3
  \item target\_recall@5
  \item target\_recall@10
\end{itemize}

sowie Family-Recall-Metriken verwendet.

\subsection{Discounted Cumulative Gain}

Järvelin und Kekäläinen führten DCG und nDCG zur Bewertung gewichteter Ranglisten ein \cite{jarvelin2002dcg}.

\begin{align}
  DCG@K = \sum_{i=1}^{K} \frac{2^{rel_i}-1} {\log_2(i+1)}
\end{align}

mit

\begin{description}
  \item[$rel_i$] Relevanz am Rangplatz $i$
\end{description}

\subsection{Normalized Discounted Cumulative Gain}

Die Normalisierung erfolgt mittels

\begin{align}
  nDCG@K = \frac{DCG@K}{IDCG@K}
\end{align}

Der Wertebereich liegt zwischen 0 und 1.

\subsection{Hierarchical nDCG}

Für technische Katalogsysteme besitzen Fehler unterschiedliche Schweregrade.
Daher wird eine hierarchische Relevanzfunktion verwendet:

\begin{align}
  rel= \begin{cases} 3 & \text{exaktes Ziel}\\ 2 & \text{gleiche Familie}\\ 0 & \text{sonst} \end{cases}
\end{align}

Diese Erweiterung erlaubt die differenzierte Bewertung von Variantenfehlern innerhalb derselben Produktfamilie.

\subsection{Candidate Coverage}

Die Candidate Coverage bewertet die maximale Leistungsfähigkeit des Retrievalsystems vor dem eigentlichen Ranking.
Hierbei wird gemessen, ob sich das korrekte Ziel überhaupt im Kandidatenpool befindet.

\subsection{Data Leakage Detection}

Kaufman et al. definieren Data Leakage als unbeabsichtigte Nutzung von Informationen, die zum Zeitpunkt der Vorhersage nicht verfügbar sein dürfen \cite{kaufman2012leakage}.
Die Verhinderung von Leakage stellt eine zentrale Voraussetzung für valide Evaluierungen und reproduzierbare Modellqualität dar.
